
\documentclass[12pt,a4paper]{article}
\usepackage{amsmath,amsthm,amssymb,mathtools}
\usepackage{tikz}
\usetikzlibrary{arrows.meta,positioning,shapes.geometric,calc,fit,backgrounds}
\usepackage{graphicx}
\usepackage{booktabs,tabularx,multirow}
\usepackage[ruled,vlined,linesnumbered]{algorithm2e}
\usepackage[most]{tcolorbox}
\usepackage{hyperref}
\usepackage[capitalize,noabbrev]{cleveref}
\usepackage[margin=1in]{geometry}
\usepackage{fancyhdr}
\usepackage{enumitem}
\usepackage{xcolor}

% ── tcolorbox environments ──────────────────────────────────────────────────
\newtcolorbox{keyinsight}[1][]{colback=blue!5!white,colframe=blue!75!black,fonttitle=\bfseries,title=Key Insight,#1}
\newtcolorbox{example}[1][]{colback=green!5!white,colframe=green!75!black,fonttitle=\bfseries,title=Example,#1}
\newtcolorbox{warningbox}[1][]{colback=red!5!white,colframe=red!75!black,fonttitle=\bfseries,title=Warning,#1}

% ── amsthm environments ─────────────────────────────────────────────────────
\newtheorem{theorem}{Theorem}[section]
\newtheorem{lemma}[theorem]{Lemma}
\newtheorem{definition}{Definition}[section]
\newtheorem{remark}{Remark}[section]

% ── Page style ───────────────────────────────────────────────────────────────
\pagestyle{fancy}
\fancyhf{}
\rhead{Session 16: RNN Introduction \& Sequence Modelling}
\lhead{Deep Learning Course -- IIITDwd}
\cfoot{\thepage}

\hypersetup{
  colorlinks=true,
  linkcolor=blue!70!black,
  urlcolor=blue!70!black,
  citecolor=green!50!black
}

% ── Convenience macros ───────────────────────────────────────────────────────
\newcommand{\xt}{x_t}
\newcommand{\ht}{h_t}
\newcommand{\yt}{y_t}
\newcommand{\Whh}{W_{hh}}
\newcommand{\Wxh}{W_{xh}}
\newcommand{\Why}{W_{hy}}

% ────────────────────────────────────────────────────────────────────────────
\begin{document}

% ============================================================
%  TITLE PAGE
% ============================================================
\begin{titlepage}
  \centering
  \vspace*{2cm}

  {\LARGE\bfseries Deep Learning Course\\[0.4em]
   Indian Institute of Information Technology, Dharwad}

  \vspace{1.5cm}
  \rule{\linewidth}{0.8pt}\\[0.4em]
  {\Huge\bfseries Session 16}\\[0.5em]
  {\Large\bfseries Recurrent Neural Networks:\\[0.3em]
   Motivation and Sequence Modelling}
  \\[0.4em]\rule{\linewidth}{0.8pt}

  \vspace{1.5cm}

  \begin{tabular}{ll}
    \textbf{Session Number:}  & 16 \\[4pt]
    \textbf{Date:}            & January 11--12, 2026 \\[4pt]
    \textbf{Instructors:}     & Prof.\ S.\ R.\ Mahadeva Prasanna \\
                              & Dr.\ Sunil Saumya \\[4pt]
    \textbf{Slide Decks:}     & \texttt{05-DL-RNN} (slides 1--63) \\
                              & \texttt{07-PyTorchDL-TransferLearning} (slides 1--29) \\[4pt]
    \textbf{Duration:}        & $\approx$ 1 hr 37 min \\[4pt]
    \textbf{Topics Covered:}  & Sequence modelling motivation; \\
                              & Limitations of feedforward networks; \\
                              & FFNN $\to$ RNN architectural transition; \\
                              & Variable-length input problem; \\
                              & Order/position blindness problem \\
  \end{tabular}

  \vfill
  {\small These notes are compiled from the lecture transcript and visual
   extractions for archival and study purposes.}
\end{titlepage}

% ============================================================
%  TABLE OF CONTENTS
% ============================================================
\tableofcontents
\newpage

% ============================================================
\section{Introduction}
\label{sec:introduction}
% ============================================================

Session~16 marks the transition in the deep learning course from spatial
models (convolutional neural networks, CNNs) to \emph{sequential} models.
Having established CNN architectures and transfer learning in the preceding
sessions, the instructors now motivate a fundamentally different class of
architectures: \textbf{Recurrent Neural Networks (RNNs)} and their
variants---Gated Recurrent Units (GRUs) and Long Short-Term Memory networks
(LSTMs).

The central question posed at the outset is: \emph{Why do feedforward neural
networks (FFNNs), despite their universal approximation capability, fail for
sequential data?} The session answers this question through two distinct
lines of argument:

\begin{enumerate}[leftmargin=2em]
  \item \textbf{Fixed window / long-range dependency failure.} A feedforward
        network operating on a sliding context window is structurally incapable
        of capturing dependencies that span more tokens than the window width.
  \item \textbf{Parallel input / order blindness.} When an entire sequence is
        flattened and fed as a single input vector, the network has no
        mechanism to distinguish word order; permuting the input produces the
        same (or nearly the same) output.
\end{enumerate}

The second half of the session, drawing from the
\texttt{07-PyTorchDL-TransferLearning} deck, formalises these problems in
the concrete setting of text classification with one-hot word vectors,
culminating in the realisation that the input dimension $N \times V$ varies
per sentence---an impossibility for a standard fully connected layer.

\begin{keyinsight}
RNNs were proposed alongside CNNs during the era of shallow neural networks,
but their true effectiveness was only demonstrated with the computational
power and training techniques of the deep learning era.  Both CNN and RNN
are \emph{structured} extensions of the basic MLP: CNNs exploit spatial
locality; RNNs exploit temporal/sequential ordering.
\end{keyinsight}

% ============================================================
\section{Background and Prerequisites}
\label{sec:background}
% ============================================================

\subsection{Data Taxonomy: Sequential vs.\ Non-Sequential}
\label{subsec:data_taxonomy}

Before introducing RNNs, the instructors establish a clear dichotomy between
the two data types that have been encountered so far in the course.

\begin{definition}[Non-Sequential Data]
\label{def:nonseq}
Data is \emph{non-sequential} when the order in which features (or rows) are
presented to a model carries no semantic information.  Permuting features
does not change the classification or regression target.
\end{definition}

\begin{definition}[Sequential Data]
\label{def:seq}
Data is \emph{sequential} when the temporal or positional ordering of
observations encodes information that is necessary to recover the target
label or prediction.  Permuting the observations destroys this information.
\end{definition}

\Cref{tab:data_types} gives concrete examples of both categories encountered
in the lecture.

\begin{table}[ht]
  \centering
  \caption{Non-sequential vs.\ sequential data examples discussed in Session~16.}
  \label{tab:data_types}
  \begin{tabularx}{\linewidth}{lXX}
    \toprule
    \textbf{Type} & \textbf{Example} & \textbf{Key Property} \\
    \midrule
    Non-sequential & Student placement data (CGPA, IQ, placement) &
      Shuffling feature columns does not change prediction \\
    Non-sequential & MNIST digit images &
      Fixed spatial grid; no temporal component \\
    Sequential & Stock price time series (open, high, close by date) &
      Shuffling rows destroys temporal patterns \\
    Sequential & Natural language text &
      Word order encodes meaning \\
    Sequential & Speech signal &
      Phoneme sequence encodes words \\
    Sequential & Video frames &
      Frame order encodes actions and events \\
    \bottomrule
  \end{tabularx}
\end{table}

\subsection{One-Hot Encoding Recap}
\label{subsec:onehot}

Throughout the session, words are represented as \emph{one-hot vectors}.
Given a vocabulary $\mathcal{V}$ of size $V$, each word $w_i$ is mapped to a
vector $\mathbf{e}_i \in \{0,1\}^{V}$ with exactly one non-zero component:

\begin{equation}
  \label{eq:onehot}
  [\mathbf{e}_i]_j =
  \begin{cases}
    1 & \text{if } j = i, \\
    0 & \text{otherwise.}
  \end{cases}
\end{equation}

A sentence of $N$ words is then a sequence
$(\mathbf{e}_{w_1}, \mathbf{e}_{w_2}, \ldots, \mathbf{e}_{w_N})$, each
vector of dimension $V$.

\subsection{Standard FFNN Review}
\label{subsec:ffnn_review}

A standard feedforward neural network with input dimension $d_{\text{in}}$,
one hidden layer of width $H$, and output dimension $d_{\text{out}}$ computes:

\begin{align}
  \label{eq:ffnn_hidden}
  \mathbf{h} &= \sigma\!\left(W_1 \mathbf{x} + \mathbf{b}_1\right), \\
  \label{eq:ffnn_output}
  \hat{\mathbf{y}} &= \text{softmax}\!\left(W_2 \mathbf{h} + \mathbf{b}_2\right),
\end{align}

where $W_1 \in \mathbb{R}^{H \times d_{\text{in}}}$,
$W_2 \in \mathbb{R}^{d_{\text{out}} \times H}$, and $\sigma$ is an
element-wise nonlinearity (e.g., ReLU or tanh).  The architecture is
\emph{fixed}: $d_{\text{in}}$ is a compile-time hyperparameter that cannot
change between inference calls.

% ============================================================
\section{Motivation for Sequence Modelling}
\label{sec:motivation}
% ============================================================

\subsection{The Significance of Sequential Order}
\label{subsec:order_significance}

Prof.\ Prasanna opens the technical discussion with a compelling linguistic
argument.  Consider the Hindi words \emph{kama} and \emph{kamala}: they
share identical constituent phonemes (\emph{ka}, \emph{ma}, \emph{la}), yet
their meanings differ entirely because of the \emph{order} in which the
phonemes appear.  A model based purely on the distribution of phoneme
occurrences would find the two words indistinguishable.

\begin{example}
The English word ``book'' illustrates context-dependent meaning driven by
sequence:
\begin{itemize}[nosep]
  \item ``I want to \textbf{book} a flight.'' \quad (\textit{book} = verb,
        to reserve)
  \item ``I want to read a \textbf{book}.'' \quad (\textit{book} = noun,
        a text)
\end{itemize}
The meaning is disambiguated entirely by the neighbouring words in the
sequence.
\end{example}

\begin{example}[title={Ball Trajectory Prediction}]
Given only the \emph{current position} of a ball on a field, all directions
of motion are equally probable.  However, given the sequence of positions
at $t{=}1,2,3,4$, the direction at $t{=}5$ can be predicted with high
confidence.  This illustrates that \emph{past sequence context} provides
the necessary information.
\end{example}

\subsection{Requirements for a Sequence Modelling Architecture}
\label{subsec:requirements}

The instructors enumerate the wishlist for any architecture designed to
handle sequential data (slide V01, \texttt{05-DL-RNN}, slide 9/63):

\begin{enumerate}[label=\textbf{R\arabic*.}, leftmargin=3em]
  \item \label{req:order} \textbf{Preserve order.} The model must process
        tokens in the order they appear; $w_t$ must be processed after
        $w_{t-1}$.
  \item \label{req:varlen} \textbf{Handle variable sequence lengths.} For the
        same class of information, different instances may have different
        sequence lengths.  The architecture must not require a fixed length.
  \item \label{req:params} \textbf{Parameter count independent of sequence
        length.} Parameters should not grow linearly (or worse) with
        sequence length; they must be \emph{shared} across time steps.
  \item \label{req:longterm} \textbf{Model long-range dependencies.}
        Information from many steps in the past must remain accessible when
        generating an output at the current step.
\end{enumerate}

\begin{keyinsight}
Requirements \ref{req:params} and \ref{req:longterm} are in tension: a model
with shared parameters that still retains long-range information requires a
form of \emph{memory} that persists and evolves across time steps.  This is
precisely the role of the hidden state $h_t$ in an RNN.
\end{keyinsight}

\subsection{Applications of Sequence Modelling}
\label{subsec:applications}

\Cref{tab:applications} summarises the sequence modelling tasks discussed in
the lecture.

\begin{table}[ht]
  \centering
  \caption{Representative sequence modelling tasks and their input/output
           structure.}
  \label{tab:applications}
  \begin{tabularx}{\linewidth}{lXX}
    \toprule
    \textbf{Task} & \textbf{Input $X$} & \textbf{Output $Y$} \\
    \midrule
    Speech recognition &
      Sequence of acoustic feature frames &
      Sequence of recognised words \\
    Machine translation &
      Sentence in source language &
      Sentence in target language \\
    Text summarisation &
      Long document ($\sim$200 words) &
      Compressed summary ($\sim$25 words) \\
    Sentiment classification &
      Text review &
      Sentiment label (positive/negative) \\
    Named entity recognition &
      Token sequence &
      Per-token entity label \\
    Video action recognition &
      Sequence of video frames &
      Action label \\
    Stock price forecasting &
      Time-ordered OHLC prices &
      Next-day closing price \\
    \bottomrule
  \end{tabularx}
\end{table}

% ============================================================
\section{Limitations of Feedforward Networks for Sequential Data}
\label{sec:ffnn_limitations}
% ============================================================

\subsection{Limitation 1: Fixed Window Cannot Model Long-Range Dependencies}
\label{subsec:fixed_window}

One naive approach to sequence modelling with FFNNs is to use a \emph{fixed
sliding window} of $k$ preceding tokens as context.  Given the current
position $t$, the network receives
$(w_{t-k}, w_{t-k+1}, \ldots, w_{t-1})$ and predicts $w_t$.

\begin{warningbox}
A fixed window of size $k$ can only capture dependencies within a span of
$k$ tokens.  Any dependency of length $> k$ is structurally invisible to
the model, regardless of how many parameters are trained.
\end{warningbox}

The lecture illustrates this with the following sentence (from Ava Soleimany,
MIT 6.S191, 2020):

\begin{quote}
``France is where I grew up, but I now live in Boston.
 I speak fluent \underline{\hspace{1.5cm}}.''
\end{quote}

The blank depends on the word ``France'' which appears many tokens earlier.
A short fixed window sees only ``\ldots live in Boston. I speak fluent'' and
might incorrectly predict ``English'' rather than ``French.''

\vspace{0.5em}
The fundamental limitation can be expressed as:

\begin{equation}
  \label{eq:fixed_window}
  P(w_t \mid w_{t-1}, w_{t-2}, \ldots, w_1)
  \;\neq\;
  P(w_t \mid w_{t-1}, w_{t-2}, \ldots, w_{t-k})
  \quad \text{when the dependency length} > k.
\end{equation}

\subsection{Limitation 2: Variable-Length Input Dimension}
\label{subsec:variable_length}

Consider a corpus $\mathcal{C}$ of sentences, each sentence
$s_i = (w_1^{(i)}, w_2^{(i)}, \ldots, w_{N_i}^{(i)})$ having $N_i$ words,
where $N_i$ varies across sentences.  Using one-hot encoding with vocabulary
size $V$, the flattened input for sentence $s_i$ has dimension $N_i \times V$.

\Cref{tab:varlen} reproduces the slide from the
\texttt{07-PyTorchDL-TransferLearning} deck (slide 7/29).

\begin{table}[ht]
  \centering
  \caption{Text input dimension as a function of sentence length.
           $V = $ vocabulary size.}
  \label{tab:varlen}
  \begin{tabular}{lcl}
    \toprule
    \textbf{Input Sentence} & \textbf{\# Words} ($N$) &
    \textbf{Network Input Size} \\
    \midrule
    We study ML               & 3 & $3 \times V$ \\
    We are in LLM class       & 5 & $5 \times V$ \\
    Today we study neural networks & 5 & $5 \times V$ \\
    Deep learning is powerful & 4 & $4 \times V$ \\
    \bottomrule
  \end{tabular}
\end{table}

An ANN requires a \emph{fixed} input dimension $d_{\text{in}}$.  Thus the
input dimension must be set to $d_{\text{in}} = N_{\max} \times V$, where
$N_{\max}$ is the length of the longest sentence in the corpus.  Shorter
sentences are \emph{zero-padded} to length $N_{\max}$.

\subsubsection*{Numerical Scale of the Problem}

Suppose $V = 1000$ unique words and $N_{\max} = 100$ words.  Then:

\begin{align}
  d_{\text{in}} &= N_{\max} \times V = 100 \times 1000 = 100{,}000, \\
  \#\text{weights (input}\to\text{hidden)} &= d_{\text{in}} \times H
    = 100{,}000 \times H.
\end{align}

For even a modest hidden layer of $H = 3$, this yields $3 \times 10^5$
weights---applied \emph{identically} to a sentence with only 5 words (which
contributes a mere $5 \times 1000 = 5{,}000$ non-zero inputs, with the
remaining $95{,}000$ inputs being zero padding).

\begin{keyinsight}
Zero-padding wastes computation proportionally to $(N_{\max} - N_i)/N_{\max}$
for each sentence $s_i$.  In the worst case (a one-word sentence against a
100-word maximum) 99\% of the input is padding noise, yet 99\% of the
weight multiplications are still performed.
\end{keyinsight}

Furthermore, at test time, if a sentence arrives with $N > N_{\max}$ words,
the network has no architectural mechanism to handle it---a hard failure.

\subsection{Limitation 3: Loss of Sequence Order Information}
\label{subsec:order_loss}

Even if the variable-length problem is ``solved'' by zero-padding, the FFNN
still suffers from order blindness.  When the entire sentence is flattened
into a single input vector and fed to the network in one forward pass, the
network has no way to distinguish \emph{which word came first}.

\begin{example}[title={Order Blindness in ANN}]
Consider the sentence ``We are in LLM class'' with vocabulary
$\mathcal{V} = \{\text{We, are, in, LLM, class}\}$.  The network input is
a $25$-dimensional vector (5 words $\times$ 5-dimensional one-hot each),
shown schematically below:

\vspace{0.3em}
\begin{center}
\begin{tikzpicture}[
    word/.style={draw,rounded corners,fill=green!10,minimum width=1.5cm,
                 minimum height=0.6cm,font=\small},
    arr/.style={-{Stealth},thick}
  ]
  \node[word] (w1) at (0,0)   {We};
  \node[word] (w2) at (1.8,0) {are};
  \node[word] (w3) at (3.6,0) {in};
  \node[word] (w4) at (5.4,0) {LLM};
  \node[word] (w5) at (7.2,0) {class};

  \node[draw,fill=yellow!20,minimum width=3cm,minimum height=0.7cm,
        font=\small] (flat) at (3.6,-1.5) {Flattened: 25-dim vector};

  \draw[arr] (w1.south) -- ++(0,-0.5) -| (flat.north west);
  \draw[arr] (w2.south) -- ++(0,-0.5) -| (flat.north);
  \draw[arr] (w3.south) -- (flat.north);
  \draw[arr] (w4.south) -- ++(0,-0.5) -| (flat.north east);
  \draw[arr] (w5.south) -- ++(0,-0.5) -| (flat.east |- flat.north);

  \node[draw,fill=orange!15,minimum width=1.4cm,minimum height=0.6cm,
        font=\small] (h) at (3.6,-3.0) {Hidden: 3};
  \node[draw,fill=red!15,minimum width=1.2cm,minimum height=0.6cm,
        font=\small] (o) at (3.6,-4.2) {Output: 1};

  \draw[arr] (flat.south) -- (h.north);
  \draw[arr] (h.south) -- (o.north);

  \node[font=\footnotesize,text=gray] at (7.5,-1.5)
    {All words fed simultaneously};
  \node[font=\footnotesize,text=gray] at (7.5,-2.0)
    {$\Rightarrow$ order is invisible};
\end{tikzpicture}
\end{center}
\vspace{0.3em}
The sentence ``class LLM in are We'' produces an \emph{identical} input
vector (up to reordering within the flattened representation), so the
network's output is the same regardless of word order.
\end{example}

\begin{remark}
\label{rem:order_loss}
This is a fundamental architectural limitation, not a data or training
deficiency.  No amount of additional data or regularisation can teach an
ANN to distinguish word order when the order information is discarded before
the first matrix multiplication.
\end{remark}

\subsection{Summary of FFNN Limitations for Sequential Data}
\label{subsec:ffnn_summary}

\begin{table}[ht]
  \centering
  \caption{Summary of feedforward network limitations for sequential data.}
  \label{tab:ffnn_limits}
  \begin{tabularx}{\linewidth}{lXX}
    \toprule
    \textbf{Limitation} & \textbf{Description} & \textbf{Consequence} \\
    \midrule
    Fixed input size &
      FFNN requires all inputs to have identical dimension &
      Cannot process variable-length sequences without padding \\
    Long-range dependency &
      Fixed window of size $k$ misses dependencies $> k$ &
      Important contextual cues (e.g., ``France $\to$ French'') are lost \\
    Order blindness &
      Parallel input processing ignores token position &
      Permutations of the same words produce the same output \\
    Parameter explosion &
      Zero-padding inflates input dimension to $N_{\max} \times V$ &
      Computation/memory scale with worst-case sequence length \\
    Poor generalisation &
      Model trained on max length $N_{\max}$ fails on $N > N_{\max}$ &
      Hard failure at test time on unseen longer sequences \\
    No parameter sharing &
      Each position has independent weight matrix in a position-wise MLP &
      The same word at different positions is processed differently \\
    \bottomrule
  \end{tabularx}
\end{table}

% ============================================================
\section{From FFNN to RNN: The Architectural Transition}
\label{sec:ffnn_to_rnn}
% ============================================================

\subsection{Conceptual Transition}
\label{subsec:conceptual_transition}

The core idea behind the FFNN $\to$ RNN transition is to \emph{unfold} the
network across time and then introduce \emph{weight sharing} across positions,
together with a \emph{recurrent connection} that carries information from
the previous time step.

Prof.\ Prasanna explains the progression in three steps:

\begin{enumerate}[leftmargin=2em]
  \item \textbf{Multiple FFNNs side by side.} Place one FFNN per token
        position, processing each token independently.  This is conceptually
        correct for position-by-position processing but does not share
        parameters and grows with sequence length.
  \item \textbf{Collapse to simplified notation.} Represent each FFNN
        (input layer $\to$ hidden layers $\to$ output layer) as a single
        compact symbol $x \to [h] \to y$.
  \item \textbf{Add the recurrent connection.} Instead of replicating the
        compact symbol horizontally (one per position), \emph{rotate it} and
        add a self-loop (arrow $C$) on the hidden node.  This loop carries the
        hidden state $h_{t-1}$ forward to time $t$, sharing the same weights
        across all time steps.
\end{enumerate}

\subsection{FFNN vs.\ RNN Architecture Diagram}
\label{subsec:arch_diagram}

\Cref{fig:ffnn_vs_rnn} reproduces the key comparison from slide V03
(\texttt{05-DL-RNN}, slide 15/63).

\begin{figure}[ht]
  \centering
  \begin{tikzpicture}[
      node distance=1.6cm,
      layer/.style={draw,rounded corners,fill=blue!10,minimum width=1.2cm,
                    minimum height=0.8cm,font=\small\bfseries},
      hlayer/.style={draw,rounded corners,fill=yellow!20,minimum width=1.2cm,
                     minimum height=0.8cm,font=\small\bfseries},
      olayer/.style={draw,rounded corners,fill=red!15,minimum width=1.2cm,
                     minimum height=0.8cm,font=\small\bfseries},
      arr/.style={-{Stealth[length=6pt]},thick},
      label/.style={font=\small\itshape}
    ]

    %% ─── FFNN (top) ──────────────────────────────────────────────────
    \node[font=\bfseries\small] at (-3.5,3.6) {Feedforward Neural Network};
    \node[layer]  (fx) at (-3.5,2.6) {$x$};
    \node[hlayer] (fh) at (-3.5,1.3) {$h$};
    \node[olayer] (fy) at (-3.5,0.0) {$\hat{y}$};

    \draw[arr] (fx) -- (fh) node[midway,right,label]{$W_1$};
    \draw[arr] (fh) -- (fy) node[midway,right,label]{$W_2$};

    \node[font=\footnotesize,text=gray] at (-3.5,-0.7)
      {(simplified: one block per layer)};

    %% ─── RNN (bottom) ────────────────────────────────────────────────
    \node[font=\bfseries\small] at (2.5,3.6) {Recurrent Neural Network};

    % Unrolled version
    \foreach \t/\lbl in {0/t{-}1,2.5/t,5.0/t{+}1}{
      \node[layer]  (rx\t) at (\t+0.0, 2.6) {$x_{\lbl}$};
      \node[hlayer] (rh\t) at (\t+0.0, 1.3) {$h_{\lbl}$};
      \node[olayer] (ry\t) at (\t+0.0, 0.0) {$y_{\lbl}$};

      \draw[arr] (rx\t) -- (rh\t);
      \draw[arr] (rh\t) -- (ry\t);
    }

    % Recurrent connections
    \draw[arr] (rh0) -- (rh2.5) node[midway,above,label]{$\Whh$};
    \draw[arr] (rh2.5) -- (rh5.0) node[midway,above,label]{$\Whh$};

    % Ellipses
    \node at (-1.0,1.3) {$\cdots$};
    \node at (6.6,1.3)  {$\cdots$};

    %% ─── Compact RNN symbol ──────────────────────────────────────────
    \node[font=\bfseries\small] at (10.0,3.6) {Compact form};
    \node[layer]  (cx) at (10.0,2.6) {$x_t$};
    \node[hlayer] (ch) at (10.0,1.3) {$h_t$};
    \node[olayer] (cy) at (10.0,0.0) {$y_t$};

    \draw[arr] (cx) -- (ch) node[midway,right,label]{$B$};
    \draw[arr] (ch) -- (cy) node[midway,right,label]{$A$};

    % Self-loop
    \draw[arr] (ch.east) to[out=0,in=0,looseness=3]
               node[right,label]{$C$ ($\Whh$)} (ch.east |- ch.south)
               to[out=180,in=270] (ch.south);

    % Labels for connections
    \node[font=\footnotesize,text=blue!70!black] at (10.0,-0.8)
      {$A$: output conn.\ ($h \to y$)};
    \node[font=\footnotesize,text=blue!70!black] at (10.0,-1.2)
      {$B$: input conn.\ ($x \to h$)};
    \node[font=\footnotesize,text=blue!70!black] at (10.0,-1.6)
      {$C$: recurrent conn.\ ($h \to h$)};

  \end{tikzpicture}
  \caption{Comparison of FFNN (left) and RNN (right).  The RNN unrolled over
           time (centre) shares weights $\Whh$, $\Wxh$, $\Why$ across all
           time steps.  The compact form (right) uses a self-loop on the
           hidden node to denote the recurrence.}
  \label{fig:ffnn_vs_rnn}
\end{figure}

\subsection{RNN Forward Pass Equations}
\label{subsec:forward_pass}

The RNN computes, at each time step $t$:

\begin{align}
  \label{eq:rnn_hidden}
  \ht &= \tanh\!\left(\Wxh\, \xt + \Whh\, h_{t-1} + \mathbf{b}_h\right), \\
  \label{eq:rnn_output}
  \yt &= \text{softmax}\!\left(\Why\, \ht + \mathbf{b}_y\right),
\end{align}

where:
\begin{itemize}[nosep]
  \item $\xt \in \mathbb{R}^{d_x}$ is the input at time $t$,
  \item $h_{t-1} \in \mathbb{R}^{d_h}$ is the hidden state from the previous
        time step,
  \item $\Wxh \in \mathbb{R}^{d_h \times d_x}$ maps input to hidden,
  \item $\Whh \in \mathbb{R}^{d_h \times d_h}$ maps hidden to hidden
        (the recurrent weight),
  \item $\Why \in \mathbb{R}^{d_y \times d_h}$ maps hidden to output.
\end{itemize}

\begin{keyinsight}
The three weight matrices $\Wxh$, $\Whh$, $\Why$ are \textbf{shared across
all time steps}.  This is the mechanism by which the number of parameters
becomes \emph{independent} of sequence length, satisfying requirement
\ref{req:params} from \cref{subsec:requirements}.
\end{keyinsight}

\subsection{How the RNN Addresses Each FFNN Limitation}

\begin{table}[ht]
  \centering
  \caption{How RNN architecture addresses the limitations of FFNNs.}
  \label{tab:rnn_solutions}
  \begin{tabularx}{\linewidth}{lX}
    \toprule
    \textbf{FFNN Limitation} & \textbf{RNN Solution} \\
    \midrule
    Fixed input size &
      Each token $x_t$ is fed one at a time; no need to specify sequence
      length at network-design time \\
    Long-range dependency &
      $h_t$ accumulates information from all previous steps; theoretically
      unlimited context window (bounded in practice for vanilla RNN) \\
    Order blindness &
      Tokens are processed sequentially; $h_t$ at step $t$ encodes the
      entire history $x_1, \ldots, x_t$ in a fixed-size summary \\
    Parameter explosion &
      $\Wxh, \Whh, \Why$ are shared; parameter count is $O(d_h^2 + d_h d_x)$,
      independent of $N$ \\
    Poor generalisation &
      The same shared weights can process sequences of any length at test time \\
    No parameter sharing &
      Weight sharing is the \emph{defining feature} of the recurrent connection \\
    \bottomrule
  \end{tabularx}
\end{table}

% ============================================================
\section{Detailed Analysis: Text Input and FFNN Problems}
\label{sec:text_input}
% ============================================================

\subsection{Text Representation Pipeline}
\label{subsec:text_pipeline}

Dr.\ Saumya's portion of the session works through a concrete text example
to make the FFNN limitations tangible.  The pipeline for representing a
sentence as input to a neural network is:

\begin{enumerate}[leftmargin=2em]
  \item \textbf{Tokenise} the sentence into individual words.
  \item \textbf{Build a vocabulary} $\mathcal{V}$ of all unique tokens.
  \item \textbf{Assign integer indices} to each token.
  \item \textbf{One-hot encode} each token using \cref{eq:onehot}.
  \item \textbf{Flatten} all one-hot vectors into a single input vector of
        dimension $N \times V$.
\end{enumerate}

\Cref{alg:text_to_input} formalises this pipeline.

\begin{algorithm}[ht]
  \caption{Text-to-Input Vector for an ANN}
  \label{alg:text_to_input}
  \SetAlgoLined
  \KwIn{Sentence $s = (w_1, w_2, \ldots, w_N)$; vocabulary $\mathcal{V}$,
        $|\mathcal{V}| = V$; max length $N_{\max}$}
  \KwOut{Flattened input vector $\mathbf{x} \in \mathbb{R}^{N_{\max} \times V}$}

  \BlankLine
  Initialise $\mathbf{x} \leftarrow \mathbf{0}_{N_{\max} \times V}$
  \tcp*[r]{zero padding by default}

  \For{$t \leftarrow 1$ \KwTo $N$}{
    $i \leftarrow \text{index of } w_t \text{ in } \mathcal{V}$\;
    $\mathbf{e}_t \leftarrow \mathbf{0}_V$\;
    $[\mathbf{e}_t]_i \leftarrow 1$
    \tcp*[r]{one-hot vector}
    $\mathbf{x}[(t-1)V : tV] \leftarrow \mathbf{e}_t$
    \tcp*[r]{insert into flattened vector}
  }
  \tcp*[r]{positions $N+1$ to $N_{\max}$ remain zero-padded}
  \Return $\mathbf{x}$\;
\end{algorithm}

\subsection{Why Parallel Feeding Destroys Order}
\label{subsec:parallel_destroys_order}

When $\mathbf{x}$ is constructed as above and fed to a standard ANN, the
hidden layer computes:

\begin{equation}
  \label{eq:ann_hidden}
  \mathbf{h} = \sigma(W_1 \mathbf{x} + \mathbf{b}_1),
\end{equation}

where $W_1 \in \mathbb{R}^{H \times (N_{\max} V)}$.  Each column of $W_1$
is associated with a specific position-word combination, not with a word
independent of its position.  Crucially, all components of $\mathbf{x}$
enter the summation simultaneously---there is no mechanism by which the
computation at position $t$ can condition on the result of computation at
position $t-1$.

\begin{remark}
This is fundamentally different from the way humans process language.  When
reading ``We are in LLM class'', the human brain processes each word
sequentially, updating an internal representation after each word.  The ANN
instead ``sees'' all five words at once, without any notion of what came
before what.
\end{remark}

\subsection{The Stochastic Model Predecessor: Hidden Markov Models}
\label{subsec:hmm}

Prof.\ Prasanna notes that before deep learning models, \emph{Hidden Markov
Models (HMMs)} were the dominant approach to sequence modelling (e.g., in
speech recognition and NLP).  HMMs encode sequence knowledge probabilistically
via state-transition probabilities:

\begin{equation}
  \label{eq:hmm}
  P(w_1, w_2, \ldots, w_T) = \prod_{t=1}^{T} P(w_t \mid w_{t-1}),
\end{equation}

(under the first-order Markov assumption).  The success of HMMs demonstrated
that sequence knowledge could be modelled formally, inspiring the development
of recurrent neural architectures that could represent richer, learned
sequence dynamics.

% ============================================================
\section{RNN Architecture: Formal Development}
\label{sec:rnn_formal}
% ============================================================

\subsection{Unrolled Computation Graph}
\label{subsec:unrolled}

The RNN processes a sequence $(x_1, x_2, \ldots, x_T)$ by unrolling the
recurrent computation across $T$ time steps.  \Cref{fig:rnn_unrolled}
shows the unrolled computation graph.

\begin{figure}[ht]
  \centering
  \begin{tikzpicture}[
      node distance=2.2cm,
      block/.style={draw,rounded corners,fill=yellow!20,minimum width=1.0cm,
                    minimum height=0.8cm,font=\small},
      input/.style={draw,rounded corners,fill=blue!12,minimum width=0.9cm,
                    minimum height=0.7cm,font=\small},
      output/.style={draw,rounded corners,fill=red!12,minimum width=0.9cm,
                     minimum height=0.7cm,font=\small},
      arr/.style={-{Stealth[length=5pt]},thick},
      lab/.style={font=\scriptsize,midway,above}
    ]

    \node[input]  (x0) at (0.0,0) {$x_1$};
    \node[input]  (x1) at (2.2,0) {$x_2$};
    \node[input]  (x2) at (4.4,0) {$x_3$};
    \node[input]  (x3) at (6.6,0) {$\cdots$};
    \node[input]  (x4) at (8.8,0) {$x_T$};

    \node[block]  (h0) at (0.0,1.6) {$h_1$};
    \node[block]  (h1) at (2.2,1.6) {$h_2$};
    \node[block]  (h2) at (4.4,1.6) {$h_3$};
    \node         (hd) at (6.6,1.6) {$\cdots$};
    \node[block]  (h4) at (8.8,1.6) {$h_T$};

    \node[output] (y0) at (0.0,3.2) {$y_1$};
    \node[output] (y1) at (2.2,3.2) {$y_2$};
    \node[output] (y2) at (4.4,3.2) {$y_3$};
    \node[output] (y4) at (8.8,3.2) {$y_T$};

    % Input to hidden
    \draw[arr] (x0) -- (h0) node[lab,right]{$\Wxh$};
    \draw[arr] (x1) -- (h1) node[lab,right]{$\Wxh$};
    \draw[arr] (x2) -- (h2) node[lab,right]{$\Wxh$};
    \draw[arr] (x4) -- (h4) node[lab,right]{$\Wxh$};

    % Hidden to output
    \draw[arr] (h0) -- (y0) node[lab,right]{$\Why$};
    \draw[arr] (h1) -- (y1) node[lab,right]{$\Why$};
    \draw[arr] (h2) -- (y2) node[lab,right]{$\Why$};
    \draw[arr] (h4) -- (y4) node[lab,right]{$\Why$};

    % Recurrent connections
    \draw[arr] (h0) -- (h1) node[font=\scriptsize,midway,above]{$\Whh$};
    \draw[arr] (h1) -- (h2) node[font=\scriptsize,midway,above]{$\Whh$};
    \draw[arr] (h2) -- (hd) node[font=\scriptsize,midway,above]{$\Whh$};
    \draw[arr] (hd) -- (h4) node[font=\scriptsize,midway,above]{$\Whh$};

    % Initial hidden state
    \node[font=\small] (h_init) at (-1.6,1.6) {$h_0{=}\mathbf{0}$};
    \draw[arr] (h_init) -- (h0) node[font=\scriptsize,midway,above]{$\Whh$};

    % Annotation
    \draw[decorate,decoration={brace,amplitude=6pt,mirror},
          draw=blue!60,thick]
      (0.0,-0.6) -- (8.8,-0.6)
      node[midway,below=8pt,font=\footnotesize,text=blue!60]
        {same $\Wxh$, $\Whh$, $\Why$ shared across all $T$ steps};

  \end{tikzpicture}
  \caption{Unrolled RNN computation graph.  The initial hidden state $h_0$
           is typically set to zero.  At each step $t$, the hidden state
           $h_t$ is computed from the current input $x_t$ and the previous
           hidden state $h_{t-1}$ using the \emph{same} weight matrices
           $\Wxh$ and $\Whh$.}
  \label{fig:rnn_unrolled}
\end{figure}

\subsection{Hidden State as Sequence Memory}
\label{subsec:hidden_state_memory}

The hidden state $h_t$ is the core mechanism by which an RNN remembers
past context.  Substituting \cref{eq:rnn_hidden} recursively:

\begin{align}
  h_1 &= \tanh(\Wxh x_1 + \Whh h_0 + \mathbf{b}_h), \\
  h_2 &= \tanh(\Wxh x_2 + \Whh h_1 + \mathbf{b}_h), \\
      &\;\vdots \notag \\
  h_t &= \tanh\!\Bigl(\Wxh x_t + \Whh \tanh\!\bigl(\Wxh x_{t-1} + \cdots
         \bigr) + \mathbf{b}_h\Bigr).
\end{align}

Thus $h_t$ is a \emph{nonlinear function of all past inputs}
$x_1, \ldots, x_t$.  This is the sense in which the RNN has ``memory.''

\begin{keyinsight}
The memory of an RNN is not an explicit data structure---it is \emph{encoded
in the activations of the hidden state}.  The hidden state $h_t$ is a
$d_h$-dimensional summary of the entire history $x_1, \ldots, x_t$.  The
quality of this summary---and how far back in time it can reach---is
determined by $d_h$ and the spectral properties of $\Whh$.
\end{keyinsight}

\subsection{Sequence Modelling Tasks and RNN Variants}
\label{subsec:rnn_variants}

As introduced by Prof.\ Prasanna, the RNN family includes several variants
that address specific shortcomings of the vanilla architecture.
\Cref{tab:rnn_variants} provides an overview.

\begin{table}[ht]
  \centering
  \caption{Overview of the RNN family.}
  \label{tab:rnn_variants}
  \begin{tabularx}{\linewidth}{lXX}
    \toprule
    \textbf{Model} & \textbf{Key Feature} & \textbf{Primary Use Case} \\
    \midrule
    Vanilla RNN &
      Single recurrent connection; $h_t = \tanh(\Wxh x_t + \Whh h_{t-1})$ &
      Short-term sequence modelling \\
    GRU (Gated Recurrent Unit) &
      Update gate and reset gate control information flow &
      Medium-range dependencies; fewer params than LSTM \\
    LSTM (Long Short-Term Memory) &
      Separate cell state $c_t$ and hidden state $h_t$; three gates &
      Long-range dependencies; speech, text \\
    Bidirectional RNN &
      Two RNNs: left-to-right and right-to-left; outputs concatenated &
      Tasks requiring full sentence context (e.g., NER) \\
    TDNN (Time Delay Neural Network) &
      Fixed dilation pattern on input; no true recurrence &
      Acoustic modelling in speech \\
    \bottomrule
  \end{tabularx}
\end{table}

\begin{remark}
As discussed in the session, vanilla RNNs still struggle with very
long-range dependencies in practice due to the \emph{vanishing gradient
problem}: gradients of the loss with respect to early time steps decay
exponentially as they are backpropagated through many tanh nonlinearities
and the recurrent weight matrix $\Whh$.  LSTM was specifically designed to
address this by introducing a cell state that can carry gradients through
long sequences with minimal decay.
\end{remark}

% ============================================================
\section{Sequence Types and Input/Output Configurations}
\label{sec:io_configs}
% ============================================================

RNNs are flexible in their input/output structure.  The lecture implicitly
introduces several configurations through the application examples.

\begin{table}[ht]
  \centering
  \caption{RNN input/output configurations.}
  \label{tab:io_configs}
  \begin{tabularx}{\linewidth}{lllX}
    \toprule
    \textbf{Configuration} & \textbf{Input} & \textbf{Output} & \textbf{Example} \\
    \midrule
    One-to-one   & Single $x$      & Single $y$       & Standard classification \\
    One-to-many  & Single $x$      & Sequence $y_{1:T}$ & Image captioning \\
    Many-to-one  & Sequence $x_{1:T}$ & Single $y$    & Sentiment classification \\
    Many-to-many (sync) & Seq $x_{1:T}$ & Seq $y_{1:T}$ & POS tagging, NER \\
    Many-to-many (async) & Seq $x_{1:T}$ & Seq $y_{1:T'}$ & Machine translation \\
    \bottomrule
  \end{tabularx}
\end{table}

% ============================================================
\section{Summary and Conclusion}
\label{sec:conclusion}
% ============================================================

Session~16 laid the motivational foundation for the entire RNN module of the
course.  The key takeaways are:

\begin{enumerate}[leftmargin=2em,label=\textbf{\arabic*.}]
  \item \textbf{Sequential data is pervasive.} Speech, text, video, and
        financial time series all require models that can exploit the ordering
        of observations.  This is not a niche requirement---it is central to
        the majority of modern AI applications.

  \item \textbf{FFNNs fail for sequential data on three fronts.}
        \begin{itemize}[nosep]
          \item Fixed input size prevents processing variable-length sequences.
          \item Parallel input processing destroys word order.
          \item Zero-padding ``fixes'' the length problem but at enormous
                computational cost and with poor generalisation to unseen lengths.
        \end{itemize}

  \item \textbf{The RNN solves all three problems.}  By processing one token
        at a time with shared weights and a recurrent hidden state, the RNN:
        (a) handles arbitrary sequence lengths, (b) preserves order (since
        $h_t$ depends on $h_{t-1}$), and (c) keeps parameter count
        $O(d_h^2 + d_h d_x)$ independent of sequence length.

  \item \textbf{The hidden state $h_t$ is the key novelty.}  It is a
        $d_h$-dimensional running summary of the sequence up to time $t$,
        updated at every step by the same matrix $\Whh$.

  \item \textbf{Vanilla RNN is a starting point.}  GRU and LSTM extend the
        vanilla RNN with gating mechanisms to address the vanishing gradient
        problem and improve long-range memory.  These will be the subject of
        subsequent sessions.
\end{enumerate}

\begin{keyinsight}
The word ``recurrent'' in \emph{Recurrent Neural Network} refers to the
\emph{recurrence} (repetition/sharing) of the same weight matrix $\Whh$
across every position in the sequence.  This is the architectural realisation
of the principle that the same learned transformation should apply at every
time step---analogous to how convolutional filters are shared spatially in a
CNN.
\end{keyinsight}

\vspace{1.5em}
\noindent\textbf{Preview of Upcoming Sessions.}  The next session will
derive the full forward and backward passes of the vanilla RNN, introduce the
Backpropagation Through Time (BPTT) algorithm, and demonstrate the vanishing
gradient problem quantitatively.  Subsequent sessions will cover GRU, LSTM,
bidirectional RNNs, and ultimately attention mechanisms and the Transformer
architecture.

% ─── Bibliography placeholder ──────────────────────────────────────────────
\vspace{2em}
\noindent\rule{\linewidth}{0.4pt}
\begin{small}
\noindent\textbf{References and Source Materials}\\[3pt]
Slide deck \texttt{05-DL-RNN}, January 11, 2026.
Sunil Saumya \& S.\ R.\ Mahadeva Prasanna, IIITDwd.\\[2pt]
Slide deck \texttt{07-PyTorchDL-TransferLearning}, January 12, 2026.
Sunil Saumya \& S.\ R.\ Mahadeva Prasanna, IIITDwd.\\[2pt]
Ava Soleimany, ``Deep Sequence Modeling,'' MIT 6.S191 (2020):
Recurrent Neural Networks.\\[2pt]
Lecture transcript: Vimeo \texttt{https://vimeo.com/1153669483},
Session~16, IIITDwd Deep Learning Course, January 2026.
\end{small}

\end{document}
